પ્રથમ-ક્રમની ભાષાનો \emph{અર્થવિચાર} તે ભાષાની \emph{સંરચના}થી મળે છે. સંરચનાનું એક \emph{વ્યાપકક્ષેત્ર} હોય છે અને દરેક અચળ સંકેતને તેનો એક ઘટક સોંપવામાં આવે છે. ફલન સંકેતોનું અર્થઘટન વ્યાપકક્ષેત્ર પરનાં ફલનોથી અને વિધેય સંકેતોનું ત્યાંના સંબંધોથી થાય છે. ચલોના ગણમાંથી વ્યાપકક્ષેત્રમાં જતા ફલનને \emph{ચલ-નિયુક્તિ} કહે છે. \emph{સંતોષ}નો સંબંધ સંરચના, ચલ-નિયુક્તિ અને સૂત્રને જોડે છે. $\Sat{M}{!A}[s]$ની વ્યાખ્યા $!A$ની સંરચના પર આગમનથી થાય છે. $\Sat{M}{!A}[s]$નો આધાર માત્ર $!A$માં ખરેખર આવતા સંકેતોના અર્થઘટન પર છે. ખાસ કરીને $!A$માં મુક્ત ચલ ન હોય તો તેનો આધાર $s$ પર નથી. એટલે $!A$ વાક્ય હોય ત્યારે કોઈ એક, અથવા સમકક્ષ રીતે દરેક, $s$ માટે $\Sat{M}{!A}[s]$ હોય તો $\Sat{M}{!A}$ છે.

સંતોષનો સંબંધ બધા અર્થવિચારી ખ્યાલોનો આધાર છે. દરેક સંરચના વાક્ય $!A$ને સંતોષે તો તે \emph{પ્રામાણ્ય ધરાવે છે}, $\Sat{{}}{!A}$. $\Gamma$ના દરેક વાક્યને સંતોષતી બધી $\Struct{M}$ માટે $\Sat{M}{!A}$ હોય ત્યારે અને માત્ર ત્યારે $\Gamma$માંથી $!A$ \emph{ફલિત થાય છે}, $\Gamma\Entails !A$. કોઈ સંરચના $\Gamma$ના દરેક વાક્યને સંતોષે ત્યારે અને માત્ર ત્યારે $\Gamma$ \emph{સંતોષ્ય} છે, અન્યથા તે \emph{અસંતોષ્ય} છે. આ ખ્યાલો પરસ્પર જોડાયેલા છે. દાખલા તરીકે, $\Gamma\Entails !A$ ત્યારે અને માત્ર ત્યારે જ્યારે $\Gamma\cup\{\lnot !A\}$ અસંતોષ્ય હોય.
